\documentclass[11pt,a4paper]{article}
\usepackage[margin=25mm]{geometry}
\usepackage{amsmath,amssymb,amsthm,hyperref}
\setlength{\emergencystretch}{2em}
\newtheorem{lemma}{Lemma}
\newtheorem{proposition}[lemma]{Proposition}
\newcommand{\gen}[1]{\langle#1\rangle}
\title{A transitive-side support deficit at a unique crossing}
\author{Human Compression project: bounded L044 proof component}
\date{}
\begin{document}\maketitle
Let $G=\gen{\rho_0,\ldots,\rho_{r-1}}\leq A_n$ be transitive,
$n\geq12$, with a string C-marking and every maximal omission
intransitive. Suppose the unique eligible $i$-edge is
$\{a,b\}$, where $0<i<r-1$ and $A,B$ are the two $G_i$-orbits.
Assume $L=G_{<i}$ is transitive on $A$ and fixes $B$ pointwise,
and that some generator of index greater than $i$ acts
nontrivially on both orbits. These are precisely the
transitive-left-side hypotheses used below; they do not
assume an intersection property for a projected group.

\begin{lemma}[An even centralizer on paired rows]\label{centralizer}
Let $p\geq4$ be even. A group $H$ preserves $p$ disjoint
pairs and induces the full natural $S_p$ on those pairs.
Every involution centralizing $H$ on their $2p$ points is even.
\end{lemma}
\begin{proof}
There are either one or two $H$-orbits, because an orbit meets
all the pairs in a constant positive number of points.
If $H$ is transitive, a nonidentity centralizing involution
has no fixed point: commutation transports any fixed point
to its entire orbit. It is therefore a product of $p$
transpositions, which is even.
If there are two orbits, each meets every pair once. The
natural action on each orbit is $S_p$ and has trivial
centralizer. A commuting permutation either preserves both
orbits, in which case it is the identity, or exchanges them.
In the latter case a commuting involution again consists
of $p$ transpositions. This proves the claim in both cases.
\end{proof}

\begin{proposition}[Transitive-side support and strict endpoint exclusion]
The largest shared label is $h=i+1$. If $h=r-1$, then
$r\leq\lfloor(n-1)/2\rfloor$. If $h<r-1$, put
$X=\operatorname{supp}(G_{>h})$. Then
\[
 X\subseteq B\setminus\{b\},\qquad
 |\operatorname{Fix}(G_{>h})|\geq |A|+1\geq2h+1,
 \qquad h\leq\frac{n-|X|-1}{2}.
\]
At equality in the last bound, $|A|=2h$ and the fixed set is
exactly $A\cup\{b\}$.
\end{proposition}
\begin{proof}
Each generator of index greater than $i$ centralizes $L$.
Its restriction to $A$ is therefore either the identity or
fixed-point-free. But every generator of index at least $i+2$
fixes $a$, by the unique-crossing lemma. Hence only
$\tau=\rho_{i+1}$ can act nontrivially on $A$, and it does.
It is fixed-point-free there, so write $|A|=2m$.
The accepted commuting-tail rank lemma, applied in the dual
index order to $L$ and $\tau$, gives
\[
                              i\leq m-1.                 \tag{1}
\]
Its hypotheses hold literally: the generators of $L$ are
supported on $A$, its action is transitive, the earlier-in-the-dual-order
involution $\tau$ centralizes it and moves $A$, and all global
omissions are intransitive.

When $h<r-1$, the tail fixes $A$ pointwise by the preceding
centralizer argument. It also fixes $b$, since every tail
index is at least $i+2$. Thus it fixes at least $2m+1$
points, and (1) gives $2m+1\geq2i+3=2h+1$.
The equality description follows directly from these two
inequalities and inclusion of $A\cup\{b\}$ in the fixed set.

It remains to exclude rank failure when $h=r-1$.
On $B$, the whole $G_i$-action is then generated by $\tau$,
since $L$ fixes $B$. Hence $B$ consists of two points $b,b'$.
Consequently $n=2m+2$ and $r=i+2\leq m+1$.
A failure of the asserted rank bound would force
\[
                        i=m-1,\qquad r=m+1.              \tag{2}
\]
The involution $\tau$ exchanges $m$ pairs on $A$ and the
pair $B$. Its evenness forces $m$ odd; moreover $n\geq12$
gives $m\geq5$.

Consider the action of $L$ on the $m$ pairs of $\tau$ in $A$.
The proof of the accepted commuting-tail lemma shows that
\emph{every} distinguished generator of $L$ is essential
in this block action: omitting a dispensable generator,
and using $\tau$ to join its two possible point orbits,
would make the corresponding global omission transitive.
Their $m-1$ block images thus form an independent involutory
string. Whiston's equality theorem identifies the block
image as $S_m$. The accepted maximal symmetric-string
classification now makes this marking an ordinary adjacent-
transposition path: degree five is covered by its complete
finite certificate, and odd degrees at least seven by its
uniform theorem. The degree-six outer action cannot occur.
Neither step assumes projected IP.

Put $H=G_{<i-1}$, the omission of the near endpoint of $L$.
Its pair action is $S_{m-1}$ on all but the anchor pair
$\{a,a'\}$, where $a'=a^\tau$. Every member of $H$ fixes
$a$ by the unique crossing, and fixes $a'$ by commutation
with $\tau$. Its only fixed points in $A$ are these two,
since its block action moves every other pair.

The crossing involution $s=\rho_i$ has the form
\[
                         s=(a\ b)u,
\]
where $u$ is an involution on $A\setminus\{a\}$; the other
point $b'$ is fixed. This follows because $(a\ b)$ is the
only cross-orbit transposition and no other point remains
in $B$ on which $s$ could act. String commutation makes
$u$ centralize $H$ on $A$. Since $u$ fixes $a$, it also
fixes the other $H$-fixed point $a'$. On the remaining
$2(m-1)$ points, Lemma~\ref{centralizer} applies with
$p=m-1$, which is even and at least four. Thus $u$ is even.
The permutation $(a\ b)u$ is odd, contradicting $G\leq A_n$.
This excludes (2) and completes the proof.
\end{proof}

\paragraph{Scope and attribution.}
This proves the numerical support conclusion and the strict
endpoint exclusion used in Cameron--Fernandes--Leemans--Mixer,
\emph{Highest rank of a polytope for $A_n$}, arXiv:1605.09173v1,
Proposition~5.22 (physical page~31), with its standing
hypotheses explicit. The exact equality fixed-set description
above does not claim the source's complete coloured-graph
classification. The proof reuses the accepted unique-crossing
and commuting-tail lemmas and all paid L047/Whiston maximal-string
inputs, including the whole degree-five certificate. No new
computation is a premise; there is no new formal declaration,
whole-original admission or question credit. The intransitive
left-side branch of Proposition~5.23 and the remaining full-IP
square/attachment rigidity are separate obligations.
\end{document}
